\section{Step 33 Charge table setup to the work that uses it}\label{crc-s33}

Compare the measured update columns with the separately recorded initializer. The prepared-table routine has less update work in these examples, but a short file may not repay table construction. State the measurement boundary before adding costs or discussing reuse across chunks.

\subsection{Expected observation}
One newly built table plus the 256-byte update costs 21009 instructions, compared with 17709 for the bit update.

\Needspace{11\baselineskip}
\subsection{Check your understanding}
Does a smaller update column prove every complete file operation is faster?

\begin{itemize}
  \item \textbf{A} Yes, because reads are free.
  \item \textbf{B} No; setup and unmeasured costs have separate boundaries.
  \item \textbf{C} Only for empty files.
\end{itemize}
